%! Factoring Trinomials — Unit 1, Lesson 1.4 — Guided Notes
\documentclass[10pt]{article}
\usepackage{atda-article}
\usepackage{atda-boxes}
\newcommand{\TallMath}[1]{$\displaystyle #1\rule[-1.4em]{0pt}{3.2em}$}
\setlength{\workrowsep}{6pt}

\begin{document}

\pageheader{Unit 1, Lesson 1.4}{Guided Notes}
% NAMESTRIP: no \namedateperiod here. The student writes their name once, on the
% cover this component is stapled behind. See references/conventions.md.

\vspace{0.15in}

\begin{objectivebox}
\small
By the end of this lesson, I will be able to\dots
\begin{enumerate}[leftmargin=1.5em, itemsep=2pt, topsep=3pt]
  \item factor a trinomial $x^2+bx+c$ by finding two integers whose \textbf{product} is $c$ and
        whose \textbf{sum} is $b$;
  \item factor a trinomial $ax^2+bx+c$ with $a \ne 1$ by \textbf{splitting the middle term} --- the
        $ac$ method --- and then grouping;
  \item factor \textbf{completely}, GCF first, and decide when a trinomial is \textbf{prime} over
        the integers;
  \item read a factored area model to recover the \textbf{dimensions} of the figure it describes,
        and say what those dimensions mean.
\end{enumerate}
\end{objectivebox}

\vspace{0.10in}

\boxguard
\begin{vocabbox}
\small
Fill in each term as we build it together.
\termblanklong{Trinomial}
\termblanklong{Quadratic trinomial in standard form}
\termblanklong{Leading coefficient}
\termblanklong{Splitting the middle term (the $ac$ method)}
\termblanklong{Prime polynomial}
\end{vocabbox}

\vspace{0.10in}

\boxguard[12]
\begin{hookbox}
\small
\textbf{The banner with no dimensions on the order.} A sign shop takes an order for a rectangular
vinyl banner. The order lists only the area:
\[ A(x) = 6x^2+19x+10 \quad \text{square feet,} \]
where $x$ is the size setting the customer chose. The cutting table needs the two \textbf{side
lengths}, and nothing on the order gives them.

Lesson 1.3 recovered the sides of a patio because the area had a shared factor. This one does not.
Guessing pairs is possible --- how many guesses would you be willing to make?

\writelines{2}
\end{hookbox}

\vspace{0.10in}

\boxguard[24]
\begin{notesbox}{1. Trinomials with Leading Coefficient $1$ --- Multiply to $c$, Add to $b$}
\small
Start from the multiplication, exactly as Lesson 1.3 did. Distributing $(x+p)(x+q)$ gives
\[ (x+p)(x+q) = x^2 + (p+q)x + pq, \]
so the middle coefficient is the \textbf{sum} of the two numbers and the constant is their
\textbf{product}. To factor $x^2+bx+c$, run that backwards: find two integers whose product is $c$
and whose sum is $b$.

\begin{itemize}[leftmargin=1.3em, itemsep=2pt, topsep=3pt]
  \item \textbf{If $c$ is positive}, both numbers carry the same sign as $b$.
  \item \textbf{If $c$ is negative}, the numbers have \blank{2.4cm} signs, and the one with the
        larger absolute value carries the sign of $b$.
\end{itemize}

\renewcommand{\arraystretch}{1.9}
\begin{tabularx}{\linewidth}{@{}p{3.2cm} p{3.6cm} X@{}}
\rowcolor{mist}
\textbf{Trinomial} & \textbf{Two numbers: product $c$, sum $b$} & \textbf{Factored form} \\
\hline
$x^2+9x+20$   & \blank{3.0cm} & \blank{3.2cm} \\
$x^2-7x+12$   & \blank{3.0cm} & \blank{3.2cm} \\
$x^2+3x-10$   & \blank{3.0cm} & \blank{3.2cm} \\
$x^2-2x-15$   & \blank{3.0cm} & \blank{3.2cm} \\
$x^2+5x+7$    & \blank{3.0cm} & \blank{3.2cm} \\
\end{tabularx}

\vspace{6pt}
\textbf{The last row is different, and the difference matters.} There are only two integer pairs
whose product is $7$:
\begin{work}
  1 \cdot 7 &= 7, \quad \text{but } 1+7 = 8 \\
  (-1)(-7) &= 7, \quad \text{but } -1+(-7) = -8
\end{work}
The search is a \emph{finite} list, so running out of pairs is a proof, not a surrender. A
polynomial that cannot be factored over the integers is \blank{2.2cm}.
\end{notesbox}

\vspace{0.10in}

\boxguard[24]
\begin{notesbox}{2. When the Leading Coefficient Is Not $1$ --- Split the Middle Term}
\small
Last night's homework already ran this engine: a trinomial has three terms, but if you
\emph{split} the middle one you get four, and four terms group. The warm-up found the rule for
choosing the split --- the two numbers multiply to $a \cdot c$ and add to $b$.

\begin{enumerate}[leftmargin=1.4em, itemsep=2pt, topsep=3pt]
  \item \textbf{GCF first}, always.
  \item Multiply $a \cdot c$.
  \item Find two integers with product $ac$ and sum $b$.
  \item Use them to rewrite $bx$ as two terms.
  \item \textbf{Group} --- Lesson 1.3, unchanged --- and check by multiplying back.
\end{enumerate}

\textbf{Back to the hook} --- the banner, where $a=6$, $b=19$, $c=10$, so $ac=60$ and the pair is
$4$ and $15$:
\begin{work}
  6x^2+19x+10 &= 6x^2+4x+15x+10 \\
              &= 2x(3x+2) + 5(3x+2) \\
              &= (3x+2)(2x+5)
\end{work}
The cutting table reads off $(3x+2)$ feet by $(2x+5)$ feet. At $x=2$ that is $8$ ft by $9$ ft, and
$8 \cdot 9 = 72 = 6(4)+19(2)+10$ square feet.

\textbf{The order of the split never matters.} Writing $6x^2+15x+4x+10$ instead groups to
$3x(2x+5)+2(2x+5)=(2x+5)(3x+2)$ --- the same two factors in the other order.

% BOXGUARD: this box cannot fit whole on a page, and \boxguard is inert inside a
% breakable tcolorbox. The natural break falls here, carrying the complete table
% (header + three rows) to the next page — a substantial chunk, not a stub. A
% \tcbbreak placed earlier balances the halves but costs a page and strands a
% single write-line, so the natural break is kept. Mirrored in the blank and key.
\renewcommand{\arraystretch}{1.9}
\begin{tabularx}{\linewidth}{@{}p{3.2cm} p{3.6cm} X@{}}
\rowcolor{mist}
\textbf{Trinomial} & \textbf{$ac$, then the pair} & \textbf{Factored form} \\
\hline
$2x^2+7x+3$    & \blank{3.0cm} & \blank{3.2cm} \\
$3x^2-10x+8$   & \blank{3.0cm} & \blank{3.2cm} \\
$4x^2+4x-15$   & \blank{3.0cm} & \blank{3.2cm} \\
\end{tabularx}
\end{notesbox}

\vspace{0.10in}

\boxguard[20]
\begin{notesbox}{3. Factoring Completely --- The Order to Work In}
\small
The order from Lesson 1.3 gains one line today:
\begin{enumerate}[leftmargin=1.4em, itemsep=2pt, topsep=3pt]
  \item \textbf{GCF first, always.} It makes every number after it smaller --- including $ac$.
  \item \textbf{Count the terms.} Four $\Rightarrow$ group. \textbf{Three $\Rightarrow$ trinomial
        methods (today).} Two gets its own methods in Lesson 1.5.
  \item \textbf{Check every factor} for something still shared, then multiply back to confirm.
\end{enumerate}
\begin{work}
  3x^3+21x^2+30x &= 3x\left(x^2+7x+10\right) \\
                 &= 3x(x+5)(x+2)
\end{work}
That inner trinomial is last night's item 13, and pulling the $3x$ out first is what kept it that
easy. Grabbing at $x^2+7x+10$ inside a cubic is much harder than grabbing at it alone.

\textbf{A negative leading coefficient is a GCF too.} Factor out the $-1$ before anything else:
$-x^2+5x-6 = -\left(x^2-5x+6\right) = -(x-2)(x-3)$.
\end{notesbox}

\vspace{0.10in}

\boxguard[20]
\begin{practicebox}
\small
\textbf{The second banner.} The same sign shop takes a second order, again with only the area
listed: $A(x) = 8x^2+22x+15$ square feet.

\begin{enumerate}[leftmargin=1.6em, itemsep=6pt, topsep=4pt]
  \item Factor $A(x)$ to find the banner's two side lengths.
\begin{work}
  8x^2+22x+15 &= 8x^2+10x+12x+15 \\
              &= 2x(4x+5) + 3(4x+5) \\
              &= (4x+5)(2x+3)
\end{work}

  \item The customer picks $x=3$. Give both side lengths and the area, in feet, with units.

\writeline

  \item A third order is for \textbf{two identical} banners with combined area
        $A(x)=2x^2+2x-24$ square feet. Factor it completely, then say what each factor tells the
        shop.
\begin{work}
  2x^2+2x-24 &= 2\left(x^2+x-12\right) \\
             &= 2(x+4)(x-3)
\end{work}
\writelines{2}
\end{enumerate}
\end{practicebox}

\end{document}
